
\textbf{Remarque 5.3.} Soient @@M0@@ un corps alg\'ebriquement clos, et @@M1@@ une surface propre et lisse sur @@M2@@. Alors le deuxi\`eme nombre de Betti @@M3@@ peut \^etre interpr\'et\'e comme \'etant @@M4@@, o\`u @@M5@@, et @@M6@@ = la deuxi\`eme classe de Chern = la caract\'eristique d'Euler-Poincar\'e = le nombre d'intersection de la diagonale de @@M7@@ avec elle-m\^eme\,; ainsi @@M8@@ a une interpr\'etation ind\'ependante de la cohomologie. L'in\'egalit\'e @@M9@@ dans cette forme pour @@M10@@ est due (avant la cohomologie \'etale) \`a IGUSA, qui a utilis\'e des r\'esultats profonds sur la structure du groupe fondamental mod\'er\'e d'une courbe alg\'ebrique moins quelques points.

\textbf{Th\'eor\`eme 5.4.} \emph{Soient @@M11@@ un corps alg\'ebriquement clos, et @@M12@@ un @@M13@@-sch\'ema propre. Alors il existe un nombre fini de courbes ferm\'ees int\`egres @@M14@@ dans @@M15@@, ayant la propri\'et\'e suivante~: pour qu'une partie @@M16@@ du sch\'ema @@M17@@ soit quasi-compacte, il faut et il suffit que les nombres d'intersection @@M18@@ des @@M19@@ restent born\'es. De plus, on peut prendre pour @@M20@@ le rang du groupe de N\'eron-Severi @@M21@@.}

En effet, comme @@M22@@ est un sous-groupe ouvert quasi-compact de @@M23@@ (4.7), pour que @@M24@@ soit quasi-compacte il faut et il suffit que l'ensemble correspondant de points de l'espace vectoriel @@M25@@ soit fini. Mais la @@M26@@-\'equivalence est \'egale \`a l'\'equivalence num\'erique (4.6 (i) @@M27@@ (iii)), donc on peut choisir parmi les courbes dans @@M28@@ des courbes qui s'en vont sur une base de l'espace dual\,; d'o\`u l'assertion.

